\begin{figure}[htbp]
\centering
\begin{tikzpicture}[scale=.88,>=Stealth,font=\small]
\foreach \c/\titulo in {0/{Elliptique},4.6/{Parabolique},9.2/{Hyperbolique}}{
\begin{scope}[xshift=\c cm]
\draw[->,gray] (-1.65,0)--(1.7,0) node[right] {\(u\)};
\draw[->,gray] (0,-1.35)--(0,1.5) node[above] {\(v\)};
\node at (0,2.05) {\titulo};
\end{scope}}
\draw[thick] (0,0) circle[radius=1];
\draw[->,thick] (1,0) arc[start angle=0,end angle=60,radius=1];
\node at (0,-1.85) {\(u^2+v^2=1\)};
\node at (0,-2.35) {\(J^2=-I\)};
\draw[thick] (3.6,-1.2)--(3.6,1.2);
\draw[->,thick] (5.6,-1.2)--(5.6,1.2);
\node at (4.6,-1.85) {\(u^2=1\)};
\node at (4.6,-2.35) {\(J^2=0\)};
\draw[thick,domain=-.9:.9,samples=50] plot ({9.2+cosh(\x)},{sinh(\x)});
\draw[thick,domain=-.9:.9,samples=50] plot ({9.2-cosh(\x)},{sinh(\x)});
\node at (9.2,-1.85) {\(u^2-v^2=1\)};
\node at (9.2,-2.35) {\(J^2=I\)};
\end{tikzpicture}
\caption{Réalisations quadratiques du TRIT. Pour
\(J_\tau=\left(\begin{smallmatrix}0&-\tau\\1&0\end{smallmatrix}\right)\),
l’exponentielle conserve \(u^2+\tau v^2\). La dégénérescence parabolique
conserve \(u\) et transporte \(v\) linéairement. Les trois diagrammes représentent
des orbites de formes quadratiques ; leur type géométrique correspond à la relation
algébrique du générateur.}
\label{fig:trit-regimenes}
\end{figure}
